% Part: first-order-logic
% Chapter: models-theories
% Section: theories

\documentclass[../../../include/open-logic-section]{subfiles}

\begin{document}

\olfileid{fol}{mat}{the}

\olsection{أمثلة لنظريات الرتبة الأولى}

\begin{ex}
لنظرية الترتيبات الخطية الصارمة في اللغة~$\Lang L_<$
نسق البديهيات
\begin{align*}
\{\quad & \lforall[\OLId{latin-ordinary}{x}][\lnot \OLId{latin-ordinary}{x} < \OLId{latin-ordinary}{x}], \\
& \lforall[\OLId{latin-ordinary}{x}][\lforall[\OLId{latin-ordinary}{y}][((\OLId{latin-ordinary}{x} < \OLId{latin-ordinary}{y} \lor \OLId{latin-ordinary}{y} <
    \OLId{latin-ordinary}{x}) \lor \OLId{latin-ordinary}{x} = \OLId{latin-ordinary}{y})]], \\
& \lforall[\OLId{latin-ordinary}{x}][\lforall[\OLId{latin-ordinary}{y}][\lforall[\OLId{latin-ordinary}{z}][((\OLId{latin-ordinary}{x} < \OLId{latin-ordinary}{y}
      \land \OLId{latin-ordinary}{y} < \OLId{latin-ordinary}{z}) \lif \OLId{latin-ordinary}{x} < \OLId{latin-ordinary}{z})]]] \quad \}
\end{align*}
وهذه البديهيات تحدّد ال!!{structure}s المقصودة تحديدًا تامًا:
فكل نموذج لها ترتيب خطي صارم؛ وبالعكس، إذا كانت $\OLId{latin-looped}{R}$ ترتيبًا
خطيًا صارمًا على مجموعة $\OLId{latin-looped}{X}$، كانت البنية $\Struct M$ ذات
$\Domain \OLId{latin-looped}{M} = \OLId{latin-looped}{X}$ و$\Assign{<}{\OLId{latin-looped}{M}} = \OLId{latin-looped}{R}$ نموذجًا للنظرية.
\end{ex}

\begin{ex}
وأما نظرية الزمر، ففي لغة تحتوي $\Obj 1$ (!!{constant}) و$\cdot$
(!!{function} ثنائية المحل)، بديهياتها:
\begin{align*}
& \lforall[\OLId{latin-ordinary}{x}][\eq[(\OLId{latin-ordinary}{x} \cdot \Obj 1)][\OLId{latin-ordinary}{x}]]\\
& \lforall[\OLId{latin-ordinary}{x}][\lforall[\OLId{latin-ordinary}{y}][\lforall[\OLId{latin-ordinary}{z}][\eq[(\OLId{latin-ordinary}{x} \cdot (\OLId{latin-ordinary}{y} \cdot \OLId{latin-ordinary}{z}))][((\OLId{latin-ordinary}{x}
          \cdot \OLId{latin-ordinary}{y}) \cdot \OLId{latin-ordinary}{z})]]]]\\
& \lforall[\OLId{latin-ordinary}{x}][\lexists[\OLId{latin-ordinary}{y}][\eq[(\OLId{latin-ordinary}{x} \cdot \OLId{latin-ordinary}{y})][\Obj 1]]]
\end{align*}
\end{ex}

\begin{ex}
وأما حساب بيانو، فتبنيه الجمل الآتية في لغة
الحساب~$\Lang L_\OLId{latin-looped}{A}$:
\begin{align*}
& \lforall[\OLId{latin-ordinary}{x}][\lforall[\OLId{latin-ordinary}{y}][(\eq[\OLId{latin-ordinary}{x}'][\OLId{latin-ordinary}{y}'] \lif \eq[\OLId{latin-ordinary}{x}][\OLId{latin-ordinary}{y}])]]\\
& \lforall[\OLId{latin-ordinary}{x}][\eq/[\Obj 0][\OLId{latin-ordinary}{x}']]\\
& \lforall[\OLId{latin-ordinary}{x}][\eq[(\OLId{latin-ordinary}{x} + \Obj 0)][\OLId{latin-ordinary}{x}]]\\
& \lforall[\OLId{latin-ordinary}{x}][\lforall[\OLId{latin-ordinary}{y}][\eq[(\OLId{latin-ordinary}{x} + \OLId{latin-ordinary}{y}')][(\OLId{latin-ordinary}{x} + \OLId{latin-ordinary}{y})']]]\\
& \lforall[\OLId{latin-ordinary}{x}][\eq[(\OLId{latin-ordinary}{x} \times \Obj 0)][\Obj 0]]\\
& \lforall[\OLId{latin-ordinary}{x}][\lforall[\OLId{latin-ordinary}{y}][\eq[(\OLId{latin-ordinary}{x} \times \OLId{latin-ordinary}{y}')][((\OLId{latin-ordinary}{x} \times \OLId{latin-ordinary}{y}) + \OLId{latin-ordinary}{x})]]]\\
& \lforall[\OLId{latin-ordinary}{x}][\lforall[\OLId{latin-ordinary}{y}][(\OLId{latin-ordinary}{x} < \OLId{latin-ordinary}{y} \liff \lexists[\OLId{latin-ordinary}{z}][\eq[(\OLId{latin-ordinary}{z}' + \OLId{latin-ordinary}{x})][\OLId{latin-ordinary}{y}])]]]\\
\intertext{ويضاف إليها جميع الجمل التي صورتها}
& (!A(\Obj 0) \land \lforall[\OLId{latin-ordinary}{x}][(!A(\OLId{latin-ordinary}{x}) \lif !A(\OLId{latin-ordinary}{x}'))]) \lif \lforall[\OLId{latin-ordinary}{x}][!A(\OLId{latin-ordinary}{x})]
\end{align*}
والجمل ذات الصورة الأخيرة لا تتناهى، ولذلك لا يتناهى نسق
البديهيات هذا. وتلك الصورة هي \emph{مخطط الاستقراء}.
(وللمخطط، على التحقيق، تفصيل أزيد مما أظهرناه هنا.)

والبديهية الأخيرة قبل مخطط الاستقراء \emph{تعريف صريح} لـ~$<$.
\end{ex}

\begin{ex}
لنظرية المجموعات الصرفة شأن في أسس الرياضيات وفي فلسفتها.
والمجموعة الصرفة ما كانت !!{element}sها مجموعات صرفة أيضًا.
فتدخل المجموعة الخالية في هذا الوصف، ولا تدخل فيه مجموعة فيها
!!a{element} ليس مجموعة. فليس في بناء المجموعات الصرفة إلا
المجموعة الخالية وما يتألف منها، دون «العناصر الأولية»، أعني
الكائنات التي ليست هي نفسها مجموعات.

وقد يُقترح لنظرية المجموعات الصرفة نسق البديهيات الآتي:
\begin{align*}
& \lexists[\OLId{latin-ordinary}{x}][\lnot \lexists[\OLId{latin-ordinary}{y}][\OLId{latin-ordinary}{y} \in \OLId{latin-ordinary}{x}]]\\
& \lforall[\OLId{latin-ordinary}{x}][\lforall[\OLId{latin-ordinary}{y}][(\lforall[\OLId{latin-ordinary}{z}][(\OLId{latin-ordinary}{z} \in \OLId{latin-ordinary}{x} \liff \OLId{latin-ordinary}{z} \in \OLId{latin-ordinary}{y})] \lif
\eq[\OLId{latin-ordinary}{x}][\OLId{latin-ordinary}{y}])]]\\
& \lforall[\OLId{latin-ordinary}{x}][\lforall[\OLId{latin-ordinary}{y}][\lexists[\OLId{latin-ordinary}{z}][\lforall[\OLId{latin-ordinary}{u}][(\OLId{latin-ordinary}{u} \in \OLId{latin-ordinary}{z} \liff
(\eq[\OLId{latin-ordinary}{u}][\OLId{latin-ordinary}{x}] \lor \eq[\OLId{latin-ordinary}{u}][\OLId{latin-ordinary}{y}]))]]]]\\
& \lforall[\OLId{latin-ordinary}{x}][\lexists[\OLId{latin-ordinary}{y}][\lforall[\OLId{latin-ordinary}{z}][(\OLId{latin-ordinary}{z} \in \OLId{latin-ordinary}{y} \liff \lexists[\OLId{latin-ordinary}{u}][(\OLId{latin-ordinary}{z} \in
        \OLId{latin-ordinary}{u} \land \OLId{latin-ordinary}{u} \in \OLId{latin-ordinary}{x})])]]]\\
\intertext{ويضاف إليها جميع الجمل التي صورتها} &
\lexists[\OLId{latin-ordinary}{x}][\lforall[\OLId{latin-ordinary}{y}][(\OLId{latin-ordinary}{y} \in \OLId{latin-ordinary}{x} \liff !A(\OLId{latin-ordinary}{y}))]]
\end{align*}
ومؤدى البديهية الأولى وجود مجموعة لا !!{element}s لها (أي إن
$\emptyset$ موجودة)؛ وتقول الثانية إن المجموعات امتدادية؛ وتقول
الثالثة إن المجموعة $\{\OLId{latin-looped}{X}, \OLId{latin-looped}{Y}\}$ موجودة لأي مجموعتين $\OLId{latin-looped}{X}$ و$\OLId{latin-looped}{Y}$؛ وتقول
الرابعة إن المجموعة $\cup \OLId{latin-looped}{X}$ موجودة لأي مجموعة $\OLId{latin-looped}{X}$، حيث $\cup \OLId{latin-looped}{X}$ هي
اتحاد جميع عناصر $\OLId{latin-looped}{X}$.

وأما ال!!{sentence}s الأخيرة، فاسمها مجتمعة \emph{مخطط الاستيعاب
الساذج}. ومؤداها وجود المجموعة $\Setabs{\OLId{latin-ordinary}{x}}{!A(\OLId{latin-ordinary}{x})}$
لكل $!A(\OLId{latin-ordinary}{x})$؛ فيبدو هذا أول النظر بديهيًا صادقًا نافعًا، وربما
بدا مما لا غنى عنه. وإنما وُصف بالسذاجة لأنه يجعل النظرية غير
قابلة للإرضاء. فلو أخذنا $!A(\OLId{latin-ordinary}{y})$ على أنها $\lnot \OLId{latin-ordinary}{y} \in \OLId{latin-ordinary}{y}$،
لزم قبول ال!!{sentence}
\[
\lexists[\OLId{latin-ordinary}{x}][\lforall[\OLId{latin-ordinary}{y}][(\OLId{latin-ordinary}{y} \in \OLId{latin-ordinary}{x} \liff \lnot \OLId{latin-ordinary}{y} \in \OLId{latin-ordinary}{y})]]
\]
ولا توجد !!{structure} تُرضى فيها هذه ال!!{sentence}.
\end{ex}

\begin{ex}
أصل النظر في مبحث \emph{الجزء والكل} علاقة \emph{الجزء}.
ولهذه العلاقة نظريات بديهية، كما للمجموعات نظريات؛ وتختلف هذه
النظريات في تصورها للعلاقة، بل قد تتعارض.

وفي لغة مبحث الجزء والكل رمز محمول ثنائي المحل واحد~$\Obj P$،
يُراد بـ$\Atom{\Obj P}{x, y}$ أن $\OLId{latin-ordinary}{x}$ جزء من~$\OLId{latin-ordinary}{y}$.
وعلى هذا القصد نسمي !!a{structure} لهذه اللغة \emph{بنية للجزء}.
وليس كل تأويل لمحمول ثنائي المحل جديرًا بهذه التسمية؛ فلا بد
أن تستوفي $\Assign{\Obj P}{\OLId{latin-looped}{M}}$ شروطًا توافق معنى «الجزء»،
فنضع هذه الشروط بديهيات لنظرية الجزء. ومن ذلك أن علاقة الجزء
ترتيب جزئي: كل كائن جزء من نفسه، وإن كان جزءًا \emph{غير فعلي}؛
ولا يكون كل واحد من كائنين مختلفين جزءًا من الآخر؛ وما كان جزءًا
من جزء كائن فهو جزء من ذلك الكائن. فهي من هذه الجهة تشبه علاقة
«هو مجموعة جزئية من»، ولا تشبه «هو عنصر من»، إذ ليست علاقة
العنصر انعكاسية ولا متعدية.
\begin{align*}
& \lforall[\OLId{latin-ordinary}{x}][\Atom{\Obj P}{x,x}] \\
& \lforall[\OLId{latin-ordinary}{x}][\lforall[\OLId{latin-ordinary}{y}][((\Part{\OLId{latin-ordinary}{x}}{\OLId{latin-ordinary}{y}} \land \Part{\OLId{latin-ordinary}{y}}{\OLId{latin-ordinary}{x}})
      \lif \eq[\OLId{latin-ordinary}{x}][\OLId{latin-ordinary}{y}])]] \\
& \lforall[\OLId{latin-ordinary}{x}][\lforall[\OLId{latin-ordinary}{y}][\lforall[\OLId{latin-ordinary}{z}][((\Part{\OLId{latin-ordinary}{x}}{\OLId{latin-ordinary}{y}} \land
        \Part{\OLId{latin-ordinary}{y}}{\OLId{latin-ordinary}{z}}) \lif \Part{\OLId{latin-ordinary}{x}}{\OLId{latin-ordinary}{z}})]]]\\
\intertext{وفوق ذلك، لكل كائنين مجموع في مبحث الجزء والكل (أي كائن
  يكون هذان الكائنان جزأين منه، ويكون أصغريًا من هذه الجهة).}  &
\lforall[\OLId{latin-ordinary}{x}][\lforall[\OLId{latin-ordinary}{y}][\lexists[\OLId{latin-ordinary}{z}][\lforall[\OLId{latin-ordinary}{u}][(\Part{\OLId{latin-ordinary}{z}}{\OLId{latin-ordinary}{u}} \liff
        (\Part{\OLId{latin-ordinary}{x}}{\OLId{latin-ordinary}{u}} \land \Part{\OLId{latin-ordinary}{y}}{\OLId{latin-ordinary}{u}}))]]]]
\end{align*}
فهذه طائفة من أصول علاقة الجزء التي يبحثها الميتافيزيقيون.
وأما ما وراءها من المبادئ، فيعسر ضبطه وصياغته من غير تقديم
علاقات معرّفة. ومن ذلك مبدأ يعدّه أكثر المشتغلين بمبحث الجزء
والكل صحيحًا: متى كان للكائن~$\OLId{latin-ordinary}{x}$ جزء فعلي~$\OLId{latin-ordinary}{y}$، كان له
جزء~$\OLId{latin-ordinary}{z}$ لا يشترك مع~$\OLId{latin-ordinary}{y}$ في شيء من الأجزاء، ويكون اندماج
$\OLId{latin-ordinary}{y}$ و$\OLId{latin-ordinary}{z}$ هو~$\OLId{latin-ordinary}{x}$.
\end{ex}

\end{document}
